%!TEX root = ../main.tex
% =============================================================================
%  第 5 章  符号说明
%  评审在问：符号清楚吗？
%    5.1 符号表（集合／参数／决策变量／对偶变量／指标量 + 上下标约定）
%  纪律：符号表覆盖全部四问，宁可多列；正文每个公式里的符号都必须能在此查到，
%        正文不得半路冒出新符号（新增符号当场登记回本表）。
% =============================================================================

\section{符号说明}\label{chap:symbols}

\subsection{符号表与上下标约定}\label{sec:notation}

核心符号、单位与索引范围见表~\ref{tab:notation}：集合与下标、参数、决策变量、\textbf{对偶变量}
（储电水价值 \(\nu_t\)、电量平衡影子价格 \(\alpha_t\)、互补松弛量 \(z^c_t,z^d_t\)，阈值律的可读性依赖它们）
以及指标量。正文其余各处只解释局部符号，同一字母不重复定义。

\begin{table}[htbp]
  \centering
  \caption{符号说明（集合、参数、决策变量、对偶变量与指标量；“---”表示无量纲；
            正文其余各处只解释局部符号，同一字母不重复定义）}
  \label{tab:notation}
  \scriptsize
  \setlength{\tabcolsep}{3pt}
  \begin{tabularx}{\linewidth}{C{0.85cm} C{2.05cm} L{6.55cm} C{1.05cm} Y}
    \toprule
    类别 & 符号 & 含义 & 单位 & 取值/来源 \\
    \midrule
    集合 & \(t\in\mathcal T\) & 时段（10~min 粒度） & --- & \(|\mathcal T|=144\) \\
    集合 & \(d,m,\omega\) & 日期 \(d\)；预报发布/决策时刻 \(m\)；情景 \(\omega\) & --- & 334 天；\(m\in\{0,6,12,18\}\) 时；\(|\Omega|=K\)（主口径 60） \\
    参数 & \(L_{d,t},\ PV_{d,t}\) & 小区负载功率、光伏实际功率 & kW & 附件 2 \\
    参数 & \(\widehat{PV}^{(m)}_{d,h}\) & \(m\) 时发布的光伏预报功率（整点值） & kW & 附件 3 \\
    参数 & \(p_{d,t}\) & 电价 & 元/kWh & 附件 1（固定价）/ 附件 4（波动价） \\
    参数 & \(\bar p_t\) & 电价的日内分时骨架（逐时段均值） & 元/kWh & 附件 1 电价剖面（A10） \\
    参数 & \(\varepsilon_{d,t}\) & 电价相对骨架的残差 & 元/kWh & \(p_{d,t}-\bar p_t\) \\
    参数 & \(\Delta t\) & 时间步长 & h & \(1/6\) \\
    参数 & \(\eta_c,\eta_d\) & 充电效率、放电效率（单边） & --- & \(0.9/0.9\)（A2） \\
    参数 & \(\overline S^{\max}\) & 储能额定容量 & kWh & \(12\,000\)（运行区间 10\%--90\%，即 1200--10\,800） \\
    参数 & \(\overline P\) & 最大充放电功率 & kW & \(5000\) \\
    参数 & \(E^{\max}\) & 单槽充/放电量上限 & kWh & \(\overline P\Delta t=833.3333\) \\
    参数 & \(S^{init}\) & 初始储电量 & kWh & \(6000\) @ 2025-01-01 \(0{:}00\) \\
    参数 & \(\kappa,\kappa^{\pm}\) & 紧急购电电价倍数；调整超额/违约倍数 & --- & \(5\)；\(1.5/0.5\) \\
    参数 & \(\beta,\alpha\) & 机会约束的缺电概率上限；CVaR 置信水平 & --- & \(0.10\)；\(0.95\)（\calref{C-11}） \\
    参数 & \(\pi_\omega\) & 情景权重 & --- & 分层样本占比，\(\sum_\omega\pi_\omega=1\) \\
    参数 & \(\rho(\cdot,\cdot)\) & 相关系数 & --- & 用于刻画电价与净负荷的共动 \\
    决策 & \(E^{plan}_{d,t},\ E^{adj}_{d,t}\) & 计划购电量（第一阶段，与 \(\omega\) 无关）、调整购电量 & kWh & \(\ge0\) \\
    决策 & \(g_{d,t}\) & 紧急购电量（第二阶段追索） & kWh & \(\ge0\) \\
    决策 & \(c_t,d_t\) & 充电量、放电量 & kWh & \(0\le c_t,d_t\le E^{\max}\) \\
    决策 & \(s_t\) & 弃光/盈余处置电量（\calref{C-06}） & kWh & \(\ge0\) \\
    决策 & \(S_t\) & 时段末储电量 & kWh & \(\in[\underline S,\overline S]\) \\
    决策 & \(u_t,v_t\) & 调整量相对计划量的增、减分解 & kWh & \(E^{adj}_t-E^{plan}_t=u_t-v_t\)，均 \(\ge0\) \\
    对偶 & \(\nu_t\ (\mu_t)\) & 储电量递推约束的乘子，即储电的边际价值（水价值） & 元/kWh
      & LP 对偶；实测 \(\in[0.4711,1.1305]\) \\
    对偶 & \(\alpha_t\) & 电量平衡约束的影子价格（\(E_t>0\) 时等于 \(p_t\)） & 元/kWh
      & LP 对偶；实测 \(\in[0.3713,1.2561]\)，恒等式偏差 \(0.000\) \\
    对偶 & \(z^{c}_t,z^{d}_t\) & 充/放互补松弛量 & 元/kWh
      & \(z^{c}_t=\alpha_t-\eta_c\nu_t\)、\(z^{d}_t=-\alpha_t+\nu_t/\eta_d\)（阈值律，144 槽违反数 0） \\
    指标 & \(C^{base}\) & 无储能基线购电费 & 元 & \calref{C-01} \\
    指标 & \(C^{plan},C^{emg}\) & 计划购电费、紧急购电费 & 元 & \calref{C-03}、\calref{C-05} \\
    指标 & \(V^{arb}\) & 储能与弃光带来的相对节省 & 元 & \(C^{base}-C^{*}\)，节省率按 \calref{C-04} \\
    \bottomrule
  \end{tabularx}
\end{table}
